% Source-checked typeset transcription of Part II, beginning at journal p. 210.
\documentclass[11pt]{article}
\usepackage[T1]{fontenc}
\usepackage[utf8]{inputenc}
\usepackage[margin=1.05in]{geometry}
\usepackage{amsmath,amssymb}
\usepackage{mathptmx}
\setlength{\parindent}{1.5em}
\setlength{\parskip}{0pt}
\begin{document}
\thispagestyle{empty}

\begin{center}\textit{Multenions and Differential Invariants.---II.}\end{center}
\begin{center}\textsc{By Alex. McAulay, Professor of Mathematics in the University of Tasmania.}\end{center}
\begin{center}(Communicated by W. B. Hardy, Sec.R.S.\quad Received June 22, 1921.)\end{center}

\noindent\textsc{§ 11. \emph{The fundamental covariant vector linity, $\eta$.}}---The meaning of a manifold as based on a quadratic differential form is here taken for granted, and this, whether, in the ordinary terminology, all or only some of its dimensions are real. Let
\begin{equation}\left.\begin{aligned}
\iota_a&=i_a, &&(a=1,2,\ldots,c),\\
\iota_b&=i_b\sqrt{-1}, &&(b=c+1,c+2,\ldots,n),\\
\rho&=\sum_{h=1}^n x_h\iota_h.
\end{aligned}\right\}\tag{1}\end{equation}
$c$ having any given value from $1$ to $n$. In the present paper we shall deal with multenions and multenion linities based on a system of vectors $\rho$ for which all the co-ordinates $x_h$ and the similar scalars are real. In the case of relativity $n=4$, $c=3$.

\indent The fundamental quadratic form is taken to be
\begin{equation}ds^2=V_0d\rho\eta d\rho=V_0d\rho'\eta'd\rho',\tag{2}\end{equation}
the sign which is given to $ds^2$, a scalar, being convenient in applications to relativity, but inconvenient in applications to pure geometry. $\eta$, a self-conjugate vector linity, is a given function of position, $\rho$; $\rho'$, a vector, an arbitrary function of $\rho$; and since the transformation from $\eta$ to $\eta'$ is given by (2) the fundamental linity $\eta$ is said to be covariant. [$\eta$ and $ds^2$ are real, but sometimes $ds$ and sometimes $ds\sqrt{-1}$ is real; when $ds$ is real it is ``a time-like interval''; when $ds\sqrt{-1}$ is real it is a length; when the question of reality is left undecided $ds$ is an interval.] At any given position $\rho'$ may be so chosen that $\eta'd\rho'=d\rho'$, that is $\eta'=1$. This imposes that $\eta$ is without singularities; it has $n$ mutually orthogonal principal directions and all its roots are real and positive, excluding zero values. Thus $\eta^{\frac12}$ is assumed to have the same orthogonal system and its roots are taken to be the positive values of the square roots of $\eta$ so that $\eta^{\frac12}$ is real and without ambiguity.

\indent We may, and from here on shall, take our previously introduced scalar $k$ (§8) to be $|\eta^{\frac12}|$, the determinant of $\eta^{\frac12}$. [We have
\[
\eta=\chi'\eta'\chi
\]
so that
\[
|\eta|=|\eta'|\times|\chi|^2=|\eta'|\times(k/k')^2,
\]
because $k/k'$ was defined by
\[
kdb=k'db'\qquad\text{or}\qquad k/k'=|\chi|.\quad\text{]}
\]

\newpage\thispagestyle{empty}\begin{center}\textit{Multenions and Differential Invariants.}\hfill 211\end{center}
It will be seen that when, in the usual notation of relativity, it is said that with Galilean co-ordinates
\[
(g_{11},g_{22},g_{33},g_{44})=(-1,-1,-1,+1),
\]
we have not that $\eta(\iota_1,\iota_2,\iota_3,\iota_4)=(-\iota_1,-\iota_2,-\iota_3,+\iota_4)$ but instead the simpler statement that $\eta=\eta^{\frac12}=1$ or
\[
\eta(\iota_1,\iota_2,\iota_3,\iota_4)=(\iota_1,\iota_2,\iota_3,\iota_4),
\]
for it will be observed that this gives
\[
ds^2=-dx_1^2-dx_2^2-dx_3^2+dx_4^2.
\]
\indent A little care is required in expressing a linity, its conjugate, reciprocal, etc., by scalars or co-ordinates. It seems sufficient to illustrate from vector linities. When we say that $\phi$ is given by the square array of $16$ scalars
\[
\left|\begin{array}{cccc}a&b&c&d\\e&f&g&h\\l&m&n&p\\q&r&s&t\end{array}\right|
\]
we imply that the co-ordinates of $\phi\iota_1$ are found in the first \emph{column} thus
\[
\phi\iota_1=a\iota_1+e\iota_2+l\iota_3+q\iota_4
\]
and so for the other columns. In the square array of $n^2$ scalars specifying $\phi$, let $a_{xy}$ be the scalar in the $x$th row and $y$th column. Thus $\phi\iota_y=\sum_{x=1}^n a_{xy}\iota_x$, or more conveniently
\begin{equation}a_{xy}=V_0\iota_x^{-1}\phi\iota_y.\tag{3}\end{equation}
Except when $c=n$ the square array of a self-conjugate linity is not symmetrical about the principal diagonal. Let $a_{xy},a'_{xy}$ be the square arrays of $\phi$ and its conjugate $\phi'$. By the meaning of conjugate we have
\[
V_0\iota_x\phi\iota_y=V_0\iota_y\phi'\iota_x.
\]
Thus instead of $a'_{yx}=a_{xy}$ (as we have when $c=n$) we have by (3)
\begin{equation}a'_{yx}=\iota_x^2\iota_y^2a_{xy},\tag{4}\end{equation}
and in particular, when $\phi$ is self-conjugate $a_{yx}=\iota_x^2\iota_y^2a_{xy}$. If we would retain the $16$ scalars $g_{ab}$ with their usual meanings as in treatises on relativity we may use $g_{(ab)}$ for the corresponding scalars specifying $\eta$. Thus from the fact that $ds^2=\sum g_{ab}dx_adx_b$ and from (3)
\[
g_{ab}=V_0\iota_a\eta\iota_b,\qquad g_{(ab)}=V_0\iota_a^{-1}\eta\iota_b,
\]
whence
\begin{equation}g_{ab}=\iota_a^2g_{(ab)}.\tag{5}\end{equation}
More details with regard to these notations and the theory of tensors will be found below in §17.

\newpage\thispagestyle{empty}\begin{center}212\hfill Prof.\ A.\ McAulay.\end{center}
\indent \emph{Note on Notations in the Present Paper.}---There will be many exceptions below to the conventions here to be described. Especially do exceptions occur in very frequently recurring symbols. Thus none of the following receive the ``cross'' notation although they should do according to the conventions,
\[
\lambda,\ E_a,\ \bar E_a,\ C_\omega,\ C.
\]
$\lambda$ is a covariant vector, and of the other four $C_\omega,C$ are contramixt linities while $E_a,\bar E_a$ are portions of contramixt linities. The explanations of the technical terms will be made as they come to be used below.

\indent (1) Contravariant multenions will have no distinguishing mark. Covariant multenions will usually be distinguished by a ``cross'' (which, because of its position, is scarcely likely to be confused with a sign of multiplication), thus
\[
\alpha^\times,\ u^\times,\ p^\times,\ q^\times,
\]
read ``alpha cross,'' etc.

\indent (2) Covariant and comixt linities will have no distinguishing mark. Contravariant and contramixt linities will usually be distinguished by a cross.

\indent (3) Contravariant and covariant linities will be denoted by Greek letters; contramixt and comixt linities by Roman capitals.

\indent (4) $\alpha,\beta,\gamma,\alpha^\times,\beta^\times,\gamma^\times$, will generally denote ``dummy'' vectors. (Eddington uses ``dummy'' in an analogous sense for a positive integer.) Dummies are always (below, though not of necessity) treated as constants in that they are not affected by the differentiations of $\nabla$.

\indent The duty of a dummy contravariant or covariant multenion may be generally described as ``to occupy, and thereby indicate, a situation into which an actual, previously defined, contravariant or covariant multenion, may legitimately be inserted.''

\noindent\textsc{§ 12. \emph{The Invariant Forms of any Multenion Expression.}}---$\eta$ has been defined as a fundamental vector linity. It is also used for the corresponding extensive linity, a multenion linity; and similarly for $\eta^{\frac12}$. Since in (2), of §11, $ds^2=(\eta^{\frac12}d\rho)^2$, $\eta^{\frac12}d\rho$ in our manifold takes the place of the $d\rho$ of multenions in a Galilean manifold. [In a Euclidean manifold all the dimensions are real. In a Galilean manifold some are imaginary.] More generally, at a given position of the manifold, $\eta^{\frac12}q$ where $q$ is a contravariant multenion, or $\eta^{-\frac12}q^\times$ where $q^\times$ is a covariant multenion, takes the place of $q$ in a Galilean manifold. $\eta^{\frac12}q$ and $\eta^{-\frac12}q^\times$ will be spoken of as the neutral forms of $q$ and $q^\times$. Indeed, the three $q$ (contravariant), $\eta q$ (covariant) and $\eta^{\frac12}q$ (neutral), are to be regarded as only different mathematical representations of the same intrinsic property, of the manifold; just as in an ordinary dynamical system, the whole system of velocities or instead the whole system of momenta equally will serve to specify the instantaneous motion.

\newpage\thispagestyle{empty}\begin{center}\textit{Multenions and Differential Invariants.}\hfill 213\end{center}
Corresponding to any multenion expression such as $V_a\cdot pqV_b(rs)$ which has a Galilean interpretation there are always the three forms ($p,q,r,s$ now being contravariant) in our manifold (contra., co., neut.).
\[
\eta^{-\frac12}V_a\cdot[\eta^{\frac12}p\eta^{\frac12}qV_b(\eta^{\frac12}r\eta^{\frac12}s)],\qquad \eta^{\frac12}V_a\cdot[\ ],\qquad V_a\cdot[\ ].
\]
The first and second of these three can \emph{always} be expressed in such a shape that only $\eta$ and not $\eta^{\frac12}$ occurs explicitly; though the expression may be very complicated as in the example chosen. It is these final forms from which $\eta^{\frac12}$ has been explicitly banished, that are regarded as invariantive.

\indent To show this, first note as an example, that the product $pqr$ can be first modified by expressing $qr$ as a sum of terms such as $V_{a+b-2c}\cdot V_aqV_br$ and then $p\cdot qr$ can be similarly modified. We have then merely to find the invariant forms (two, namely, covariant and contravariant) of $V_{a+b-2c}uv$. [The reader is reminded that in our notation $u,v,w$, have homogeneities $a,b,c$, respectively, or have the forms $u=V_ap,\ v=V_bq,\ w=V_cr$.]

\indent Let $a$ be greater or equal to $b$ and let $u,v$ be contravariant. Then most frequently all we require to note are the two cases $c=0,c=b$:
\[
\eta^{-\frac12}V_{a+b}\eta^{\frac12}u\eta^{\frac12}v=V_{a+b}uv,\qquad \eta^{\frac12}V_{a+b}\eta^{\frac12}u\eta^{\frac12}v=V_{a+b}\eta u\eta v,
\]
\[
\eta^{-\frac12}V_{a-b}\eta^{\frac12}u\eta^{\frac12}v=V_{a-b}u\eta v,\qquad \eta^{\frac12}V_{a-b}\eta^{\frac12}u\eta^{\frac12}v=V_{a-b}\eta u v.
\]
and here of course we may have $u=\eta^{-1}u^\times$ or $v=\eta^{-1}v^\times$ in each of the forms.

\indent To deal with the general value of $c$ in $V_{a+b-2c}uv$ we use $\zeta$. Now we know that if $\tau_1,\tau_2,\ldots\tau_n$ are $n$ independent contravariant vectors at a point, then their normal reciprocals $\hat{\tau}_1,\hat{\tau}_2,\ldots\hat{\tau}_n$ are $n$ independent covariant vectors. Also we know that if $(\alpha,\beta^\times)$ is any expression---bilinear in two vectors
\[
(\zeta,\zeta)=-\sum_{a=1}^n(\tau_a,\hat{\tau}_a).
\]
Hence if we insert $\zeta,\zeta$ into two places in an invariant form, of which one is occupied by a contravariant vector and the other by a covariant vector, the resulting expression will be invariantive (meaning, ``of invariant form''). Now it is easy to see by reduction to primitive units that
\begin{equation}V_{a+b-2c}uv=[(-1)^{\frac12c(c+1)}/c!]V_{a+b-2c}\cdot V_{a-c}u\zeta^{(c)}V_{b-c}\zeta^{(c)}v,\tag{1}\end{equation}
where as usual $\zeta^{(c)}$ stands for $\zeta_1\zeta_2\ldots\zeta_c$ and $\zeta_c^{(c)}$ for $V_c\zeta^{(c)}$. Hence the contravariant form of $V_{a+b-2c}uv$ is
\begin{equation}[(-1)^{\frac12c(c+1)}/c!]V_{a+b-2c}\cdot V_{a-c}u\zeta_c^{(c)}V_{b-c}\zeta_c^{(c)}v,\tag{2}\end{equation}
from which the covariant form may be written down, and either $u^\times=\eta u$ or $v^\times=\eta v$ or both may be used in place of $u$ and $v$. In my MS. work I

\newpage\thispagestyle{empty}\begin{center}214\hfill Prof.\ A.\ McAulay.\end{center}
am in the habit of using $\zeta_K^{(c)},\zeta_Q^{(c)}$ for $K\zeta_c^{(c)},Q\zeta_c^{(c)}$. Thus (2) would be $(c!)^{-1}V_{a+b-2c}\cdot V_{a-c}u\zeta_c^{(c)}V_{b-c}\eta\zeta_K^{(c)}v$. A special case of (2) often required is that the contravariant form of $V_2\omega\omega'$ where $\omega$ and $\omega'$ are contravariant and ${}_nV_2$ is
\begin{equation}V_2(V_1\zeta\omega)(V_1\eta\zeta\omega').\tag{3}\end{equation}
or more generally, for infinitesimal rotations generally, the contravariant form of $V_a\omega\iota$ is
\begin{equation}V_a(V_1\zeta\omega)(V_{a-1}\eta\zeta\iota).\tag{4}\end{equation}
\noindent\textsc{§ 13. \emph{Normal and Incident Components in Invariant Form.}}---These components occur very frequently in relativity. Thus the energy-momentum linity is density multiplied by an incident component, at any rate in its kinetic or primitive form. Normal and incident components serve to separate: (1) vector potential from electrostatic potential; (2) momentum from energy; (3) force from power; (4) current density from charge density; (5) magnetic induction from electrostatic force; (6) magnetic field strength from displacement. One should always be prepared to recognise these resolutions.

\indent In quaternions if $\alpha$ is a given vector and $\sigma$ an arbitrary one, $\sigma$ has two components with reference to $\alpha$, namely the normal component
\[
N_0\sigma=\alpha^{-1}V\alpha\sigma=\alpha V\alpha\sigma/\alpha^2
\]
and the incident component
\[
I_0\sigma=\alpha^{-1}S\alpha\sigma=\alpha S\alpha\sigma/\alpha^2.
\]
As will be familiar to users of the quaternion method, it is these components each multiplied by $\alpha^2$, rather than the components themselves, which we shall expect to meet with frequently in analytical expressions. The two
\[
N_\alpha\sigma=N\sigma=\alpha V\alpha\sigma,\qquad I_\alpha\sigma=I\sigma=\alpha S\alpha\sigma,
\]
may be referred to as the normal and incident \emph{expressions} of $\sigma$ with reference to $\alpha$. With regard to these four linities $N_0,I_0,N,I$ we have as follows
\begin{equation}\left.\begin{aligned}N_0^2&=N_0,&I_0^2&=I_0,&N_0I_0&=0=I_0N_0,\\N_0'&=N_0,&I_0'&=I_0,&N_0+I_0&=1=N_0'+I_0'\end{aligned}\right\},\tag{1}\end{equation}
\begin{equation}\left.\begin{aligned}N^2&=\alpha^2N,&I^2&=\alpha^2I,&NI&=0=IN,\\N'&=N,&I'&=I,&N+I&=\alpha^2=N'+I'\end{aligned}\right\}.\tag{2}\end{equation}
Similarly in a Galilean manifold any multenion $q$ has two components $N_0q,I_0q$ normal and incident to a given vector $\alpha$. Let
\[
q=\sum_{c=1}^n w,\qquad\text{where}\quad w=V_cq;
\]

\newpage\thispagestyle{empty}\begin{center}\textit{Multenions and Differential Invariants.}\hfill 215\end{center}
then when the linities $Nw,N_0w$, etc., have been defined, $Nq,N_0q$ have also been defined. We have
\begin{equation}\left.\begin{aligned}Nw&=\alpha V_{c+1}\alpha w,&Iw&=\alpha V_{c-1}\alpha w,\\N_0w&=\alpha^{-1}V_{c+1}\alpha w,&I_0w&=\alpha^{-1}V_{c-1}\alpha w\end{aligned}\right\}.\tag{3}\end{equation}
[If $\epsilon$ is a vectorium of homogeneity $b$, a vector $\sigma$ has a normal component $\epsilon^{-1}V_{b+1}\epsilon\sigma$ and an incident component $\epsilon^{-1}V_{b-1}\epsilon\sigma$ with reference to $\epsilon$. With reference to $\epsilon$, $w$ has many, not merely two, components given by
\[
w=\epsilon^{-1}V_{b+c}\epsilon w+\epsilon^{-1}V_{b+c-2}\epsilon w+\cdots.
\]
The terms on the right of this equation are all, separately, but not obviously, of homogeneity $c$.]

\indent Next let $\alpha$ be a given contravariant vector at a given position $\rho$ of our manifold and $q$ an arbitrary contravariant multenion at the same position and let $w=V_cq$. Also let $\alpha_v,q_v,w_v$ be the neutral forms and $\alpha^\times,q^\times,w^\times$ the covariant forms of $\alpha,q,w$ respectively, that is to say,
\begin{equation}\left.\begin{aligned}\eta^{\frac12}\alpha&=\alpha_v=\eta^{-\frac12}\alpha^\times,\\\eta^{\frac12}w&=w_v=\eta^{-\frac12}w^\times,\\\eta^{\frac12}q&=q_v=\eta^{-\frac12}q^\times\end{aligned}\right\}.\tag{4}\end{equation}
The invariant form of $\alpha^2$ is $\alpha_v^2$ or
\begin{equation}V_0\alpha\eta\alpha=\alpha_v^2=V_0\alpha^\times\eta^{-1}\alpha^\times.\tag{5}\end{equation}
With our present meanings (3) is not invariant in form. There are four modes of choosing the invariant forms of $N$ and $I$, $N_0$ and $I_0$ and generally of any multenion linities, (1) covariant, (2) contravariant, (3) comixt, (4) contramixt. We will take number (3) of these for the basis of the notation in connection with $N$ and $I$. A comixt linity means one which acting on a covariant multenion produces a covariant multenion. For this type of linity it is obvious that the invariant forms of (3) are given by
\begin{equation}\left.\begin{aligned}
Nw^\times&=V_c\alpha V_{c+1}\eta\alpha w^\times=V_c\eta^{-1}\alpha^\times V_{c+1}\alpha^\times w^\times,\\
Iw^\times&=V_c\eta\alpha V_{c-1}\alpha w^\times=V_c\alpha^\times V_{c-1}\eta^{-1}\alpha^\times w^\times,\\
N_0w^\times&=Nw^\times/V_0\alpha\eta\alpha=Nw^\times/V_0\alpha^\times\eta^{-1}\alpha^\times,\\
I_0w^\times&=Iw^\times/V_0\alpha\eta\alpha=Iw^\times/V_0\alpha^\times\eta^{-1}\alpha^\times
\end{aligned}\right\}.\tag{6}\end{equation}
(2) comes with this notation to be modified in two ways only; $\alpha^2$ is replaced by $V_0\alpha\eta\alpha$ or $V_0\alpha^\times\eta^{-1}\alpha^\times$ and in place of $N'$ and $I'$ we have a new kind of conjugate $\eta N'\eta^{-1},\eta I'\eta^{-1}$ involving the fundamental linity $\eta$. Thus we now have
\begin{equation}\left.\begin{aligned}
N^2&=V_0\alpha\eta\alpha\cdot N,&I^2&=V_0\alpha\eta\alpha\cdot I,&NI&=0=IN,\\
\eta N'\eta^{-1}&=N,&\eta I'\eta^{-1}&=I,&N+I&=V_0\alpha\eta\alpha=\eta(N'+I')\eta^{-1}
\end{aligned}\right\}.\tag{7}\end{equation}
\newpage\thispagestyle{empty}\begin{center}216\hfill Prof.\ A.\ McAulay.\end{center}
[Note that if we put $\eta\phi'\eta^{-1}=\phi_c$ we obtain for this new conjugate the familiar properties
\[
(\phi+\psi)_c=\phi_c+\psi_c,\qquad(\phi\psi)_c=\psi_c\phi_c,\qquad(\phi_c)_c=\phi.]
\]
The types of linity we have to recognise are given in the following short list.
\begin{center}\textsc{Linity Types.}\end{center}
\begin{center}\begin{tabular}{|l|c|c|c|c|c|}\hline
Name & Co-variant. & Contra-variant. & Co-mixt. & Contra-mixt. & Neutral.\\\hline
Symbol & $\xi$ & $\xi^\times$ & $X$ & $X^\times$ & $X_v=\phi$\\
Meaning & $p^\times=\xi q$ & $p=\xi^\times q^\times$ & $p^\times=Xq^\times$ & $p=X^\times q$ & $p_v=X_vq_v$\\\hline
\end{tabular}\end{center}
Thus $\xi$ converts a $q$ (contravariant) into a $p^\times$ (covariant) and similarly for the other types.

\indent The two mixed linities convert a multenion of a given type to one of the same type; the other two convert a given type to the \emph{other} type. [In our present method the mixed linities could scarcely be more unhappily named.]
\indent Corresponding to any one of the types we have, of course, by means of $\eta$, symbols of the other types with the same intrinsic significance. Thus, corresponding to the neutral form $\phi$ we have
\begin{equation}\left.\begin{aligned}\xi&=\eta^{\frac12}\phi\eta^{\frac12},&\xi^\times&=\eta^{-\frac12}\phi\eta^{-\frac12},\\X&=\eta^{\frac12}\phi\eta^{-\frac12},&X^\times&=\eta^{-\frac12}\phi\eta^{\frac12}\end{aligned}\right\}.\tag{8}\end{equation}
From the list and from (8) the following statements, in which $\xi_1,\xi_2$ mean two linities of the type $\xi$, etc., are easily established:---
\par\indent (1) $\xi_1\xi_2$ has no invariant significance, nor has $\xi_1^\times\xi_2^\times$.
\par\indent (2) $X_1X_2$ is of type $X$ and $X_1^\times X_2^\times$ is of type $X^\times$.
\par\indent (3) Although by (1), for integral values of $c$, $\xi^c$ has in general no invariant significance; in particular when $c=-1$ it has significance. $\xi^{-1}$ is of type $\xi^\times$, and $\xi^{\times-1}$ is of type $\xi$. $X^{-1}$ is of type $X$, and $X^{\times-1}$ is of type $X^\times$.
\par\indent (4) A rational function $f(X)$ is of type $X$; and $f(X^\times)$ is of type $X^\times$; the corresponding neutral form in each case being $f(\phi)$.
\par\indent (5) $\xi'$ is of type $\xi$, $\xi^{\times\prime}$ of type $\xi^\times$; so that a similar remark applies to the self-conjugate and skew parts of $\xi$ and $\xi^\times$. $\xi'$ and $\xi^{\times\prime}$ correspond to $\phi'$.
\par\indent (6) $X'$ is of type $X^\times$, $X^{\times\prime}$ of type $X$, both corresponding to $\phi'$. Again $\eta X'\eta^{-1}$ is of type $X$ and $\eta^{-1}X^{\times\prime}\eta$ of type $X^\times$, both corresponding to $\phi'$.
\par\indent (7) The identity satisfied by $\xi$ or by $\xi^\times$ has no significance. The
\newpage\thispagestyle{empty}\begin{center}\textit{Multenions and Differential Invariants.}\hfill 217\end{center}
identity satisfied by $\phi$ is the same as the identity satisfied by each of the following
\[
X=\xi\eta^{-1}=\eta\xi^\times,\qquad X^\times=\eta^{-1}\xi=\xi^\times\eta,
\]
and all their conjugates and $\eta$-conjugates.

\indent (8) In the neutral form a finite rotation is effected by a vector linity (and its extensive) $\phi$ for which $\phi'\phi=1$. This is equivalent to
\[
\xi'\eta^{-1}\xi\eta^{-1}=\eta X'\eta^{-1}X=1.
\]
From the above it will be seen that if in a Galilean manifold we would deal with $\phi',\phi'\pm\phi$, etc., then probably in the general manifold it is best to deal with $\xi$ and $\xi^\times$; if in the Galilean we would use $\phi^2,f(\phi)$ etc., then in the general manifold the use of $X,X^\times$ is indicated.

\indent We need not write down the twenty-four alternative forms of the eight equations (6), but we may note how in terms of $N,I$ the other types appear:
\begin{equation}\left.\begin{aligned}
N,I&\text{ are of type }X,\\
N'=\eta^{-1}N\eta,\quad I'=\eta^{-1}I\eta&\text{ are of type }X^\times,\\
N\eta,\ I\eta&\text{ are of type }\xi,\\
\eta^{-1}N,\ \eta^{-1}I&\text{ are of type }\xi^\times
\end{aligned}\right\}.\tag{9}\end{equation}
After the fundamental linity $\eta$, of type $\xi$, has presented itself the other three may be regarded as arising from
\[
\xi^{-1},\qquad\xi_1\xi_2^{-1},\qquad\xi_1^{-1}\xi_2.
\]
\noindent\textsc{§ 14. \emph{Riemann's and Weyl's Manifolds.}}---Suppose $\rho,\rho+\alpha,\rho+\beta$ are for our manifold the position vectors of three infinitesimally near points, P, A, B. We assume that for such small regions the metrics are intrinsically the same as in a Galilean system, and in particular that there is a definite shift of PA along PB, and of PB along PA, each called a translation, by which if A in the first shift comes to C, B also in the second shift will come to C. With arbitrary co-ordinates we cannot of course affirm that C will have the position vector $\rho+\alpha+\beta$, but instead that it will have
\[
\rho+\alpha+\beta+d_\alpha^\tau\beta
\]
where $d_\alpha^\tau$ is to be regarded, not as a symbol of operation, but merely as the abbreviation of the words
\begin{equation}d_\alpha^\tau=\text{``increment due to translation along }\alpha\text{.''}\tag{1}\end{equation}
From these preliminaries it is clear that
\begin{equation}d_\alpha^\tau\beta=d_\beta^\tau\alpha=-E_a\beta,\tag{2}\end{equation}
where $E_a$ is a vector linity which is a function of the position vector $\rho$, and is also a function of $\alpha$, linear in the latter; also $E_a\beta$ is symmetrical in $\alpha$ and $\beta$.
\newpage\thispagestyle{empty}\begin{center}218\hfill Prof.\ A.\ McAulay.\end{center}
The word translation implies something about constancy or otherwise in interval or in ``the square measure'' $ds^2$ of equation (2) of §11. In Weyl's view we are to assume that in general a \emph{translation} is such that it is impossible for the square measure to remain unaltered when by successive steps a closed path has been traced. In such a closed path beginning and ending at a given point, there will appear a change in the square measure depending on the point and the path.

\indent The square measure then is assumed to change according to the equation
\begin{equation}d_\alpha^\tau V_0\beta\eta\beta=-V_0\beta\eta\beta\cdot V_0\lambda\alpha,\tag{3}\end{equation}
where $\alpha$ and $\beta$ are arbitrary and with the same meanings as before, and $\lambda$ is a vector, a given function of $\rho$. Thus $\alpha,\beta$ must be contravariant, and $\lambda$ covariant. $E_a\beta$ is a vector neither covariant nor contravariant, which need not surprise us, because more than one position of the manifold has been used in its description.

\indent It is clear that $\eta$ is arbitrary as to a scalar factor and $\lambda$ correspondingly arbitrary. More definitely, nothing is altered in the meanings given if (1) $\eta$ is multiplied by a positive scalar $h$, any function of position, and (2) $\lambda$ is altered to $\lambda+\nabla\log h$, as appears from equation (3).

\indent \emph{Important qualification of the last sentence.}---$\eta$ here means the original \emph{vector} linity denoted by $\eta$. Whenever we speak of altering gauge by multiplying $\eta$ by $h$ this restriction is to be understood. It is clear that $\eta w$ thereby changes, not to $h\eta w$, but to $h^c\eta w$.

\indent ``Into all expressions or formulae, which exhibit analytically true metrical properties of the manifold, $\eta$ and $\lambda$ must enter in such a manner that invariance subsists (1) when there is arbitrary transformation of co-ordinates ($\rho$ replaced by $\rho'$); and (2) when $\eta,\lambda$ are replaced by $h\eta,\lambda+\nabla\log h$, whatever positive scalar function of $\rho,h$ may be.''\footnote{Weyl, \emph{Raum, Zeit, Materie}, 3rd ed., p. 110. English translation, 1922, of 4th ed. (Brose), last sentence of p. 123.}

Since (3) is true, and also what (3) becomes both when $\beta$ is changed to $\gamma$ and when changed to $\beta+\gamma$, we get the corresponding bilinear form
\begin{equation}d_\alpha^\tau V_0\beta\eta\gamma=-V_0\beta\eta\gamma\cdot V_0\lambda\alpha,\tag{4}\end{equation}
where $\alpha,\beta,\gamma$ are all arbitrary. In the translation denoted by $d_\alpha^\tau$, $\eta$ suffers the increment $-V_0\alpha\nabla\cdot\eta$ and $\beta$ and $\gamma$ the increments $-E_a\beta,-E_a\gamma$. Hence
\[
-V_0\alpha\nabla\eta\cdot V_0\beta\eta\gamma-V_0E_a\beta\eta\gamma-V_0\beta\eta E_a\gamma=-V_0\beta\eta\gamma\cdot V_0\lambda\alpha.
\]
Put
\begin{equation}\eta E_a=\Gamma_a,\tag{5}\end{equation}
so that $\Gamma_a\beta$, like $E_a\beta$, is symmetrical in $\alpha,\beta$. Thus
\[
-V_0(\gamma\Gamma_a\beta+\beta\Gamma_a\gamma)=V_0\alpha(\nabla\eta-\lambda)\cdot V_0\beta\eta\gamma.
\]
\newpage\thispagestyle{empty}\begin{center}\textit{Multenions and Differential Invariants.}\hfill 219\end{center}
Write down the two equations obtained by cyclically changing $\alpha,\beta,\gamma$, change the signs of the second and third, and then add the three equations. We obtain
\[
2V_0\alpha\Gamma_\beta\gamma=V_0\alpha\bigl[(\nabla\eta-\lambda)V_0\beta\eta\gamma-V_0\beta(\nabla\eta-\lambda)\cdot\eta\gamma-V_0\gamma(\nabla\eta-\lambda)\cdot\eta\beta\bigr],
\]
or
\begin{equation}\Gamma_\beta\gamma=\tfrac12\bigl[-V_0\beta(\nabla\eta-\lambda)\cdot\eta\gamma-V_0\gamma(\nabla\eta-\lambda)\cdot\eta\beta+(\nabla\eta-\lambda)V_0\beta\eta\gamma\bigr].\tag{6}\end{equation}
For Riemann's manifold we have merely to put $\lambda=0$. Since frequent reference must be made to this case of Riemann's manifold, we will put for the values of $E_a$ and $\Gamma_a$ when $\lambda=0$, $\bar E_a$ and $\bar\Gamma_a$ respectively.

\indent From (6)
\begin{equation}(\Gamma_\beta+\Gamma_\beta')\gamma=-V_0\beta(\nabla\eta-\lambda)\cdot\eta\gamma,\tag{7}\end{equation}
\begin{equation}(\Gamma_\beta-\Gamma_\beta')\gamma=V_1\bigl[V_2(\nabla\eta-\lambda)\eta\beta\cdot\gamma\bigr].\tag{8}\end{equation}
Thus $\Gamma_\beta'\gamma$ is obtained from (6) by reversing the signs of the second and third terms on the right. It further follows that $\Gamma_\beta'\gamma$ regarded as a linity of $\beta$, instead of $\gamma$, is self-conjugate.

\indent This, and its converse, can be shown to be true of any such function $Y_{\alpha\beta}$ which is symmetrical in $\alpha$ and $\beta$. Thus we may put for the $\iota_c$ component of $Y_{\alpha\beta}$, $\iota_cV_0\alpha\psi_c\beta$ where $\psi_c$ is a self-conjugate vector linity, or
\begin{equation}\left.\begin{aligned}
Y_{\alpha\beta}&=\sum_{c=1}^n\iota_cV_0\beta\psi_c\alpha\\
Y_\alpha'\beta&=\sum_{c=1}^n\psi_c\alpha V_0\iota_c\beta
\end{aligned}\right\},\tag{9}\end{equation}
showing that $Y_\alpha'\beta$ is self-conjugate with respect to $\alpha$. Each $\psi_c$ involves $\tfrac12n(n+1)$ scalars, so that $Y_{\alpha\beta}$ involves $\tfrac12n^2(n+1)$ scalars; when $n=4$ this number is 40.

\indent When $Y_{\alpha\beta}$ is not symmetrical (9) still holds, $\psi_c$ now being not self-conjugate, and the number of scalars being $n^3$.

\indent We shall from this point onwards take for granted, as is not difficult to prove, that \emph{at a given point it is always possible so to choose our co-ordinates that all the $\tfrac12n^2(n+1)$ differential coefficients $(\nabla\eta,\eta)$ are zero; and independently to choose the gauge $(h)$ so that $\lambda$ is zero; therefore also to make $E_a$ zero for all values of $\alpha$.} [Brose calls $h$ the calibration ratio. Selection of a definite $h$ he calls calibration.]

\indent\textsc{§ 15. \emph{Absolute Differentiation.}}---Let $\tau$ be a contravariant vector given over an assigned path or over a region of $n$ or fewer dimensions. If the region is of fewer than $n$ dimensions, we may still assign arbitrary values over a region of $n$ dimensions containing the given region, so that we can continue to apply $\nabla$ to it.

\indent In passing from the point $\rho$ to the point $\rho+\alpha$, where $\alpha$ is infinitesimal, the
\end{document}
